\documentclass[12pt,letterpaper]{article}
\usepackage{graphicx,textcomp}
\usepackage{natbib}
\usepackage{setspace}
\usepackage{fullpage}
\usepackage{color}
\usepackage[reqno]{amsmath}
\usepackage{amsthm}
\usepackage{fancyvrb}
\usepackage{amssymb,enumerate}
\usepackage[all]{xy}
\usepackage{endnotes}
\usepackage{lscape}
\newtheorem{com}{Comment}
\usepackage{float}
\usepackage{hyperref}
\newtheorem{lem} {Lemma}
\newtheorem{prop}{Proposition}
\newtheorem{thm}{Theorem}
\newtheorem{defn}{Definition}
\newtheorem{cor}{Corollary}
\newtheorem{obs}{Observation}
\usepackage[compact]{titlesec}
\usepackage{dcolumn}
\usepackage{tikz}
\usetikzlibrary{arrows}
\usepackage{multirow}
\usepackage{xcolor}
\newcolumntype{.}{D{.}{.}{-1}}
\newcolumntype{d}[1]{D{.}{.}{#1}}
\definecolor{light-gray}{gray}{0.65}
\usepackage{url}
\usepackage{listings}
\usepackage{color}

\definecolor{codegreen}{rgb}{0,0.6,0}
\definecolor{codegray}{rgb}{0.5,0.5,0.5}
\definecolor{codepurple}{rgb}{0.58,0,0.82}
\definecolor{backcolour}{rgb}{0.95,0.95,0.92}

\lstdefinestyle{mystyle}{
	backgroundcolor=\color{backcolour},   
	commentstyle=\color{codegreen},
	keywordstyle=\color{magenta},
	numberstyle=\tiny\color{codegray},
	stringstyle=\color{codepurple},
	basicstyle=\footnotesize,
	breakatwhitespace=false,         
	breaklines=true,                 
	captionpos=b,                    
	keepspaces=true,                 
	numbers=left,                    
	numbersep=5pt,                  
	showspaces=false,                
	showstringspaces=false,
	showtabs=false,                  
	tabsize=2
}
\lstset{style=mystyle}
\newcommand{\Sref}[1]{Section~\ref{#1}}
\newtheorem{hyp}{Hypothesis}


\title{Problem Set 4}
\date{Due: November 26, 2026}
\author{Elena Gyubin Lee}


\begin{document}
	\maketitle
	\section*{Instructions}
	\begin{itemize}
		\item Please show your work! You may lose points by simply writing in the answer. If the problem requires you to execute commands in \texttt{R}, please include the code you used to get your answers. Please also include the \texttt{.R} file that contains your code. If you are not sure if work needs to be shown for a particular problem, please ask.
		\item Your homework should be submitted electronically on GitHub.
		\item This problem set is due before 23:59 on Thursday November 26, 2026. No late assignments will be accepted.
	\end{itemize}



	\vspace{.5cm}
\section*{Question 1: Economics}
\vspace{.25cm}
\noindent 	
In this question, use the \texttt{prestige} dataset in the \texttt{car} library. First, run the following commands:

\begin{verbatim}
install.packages(car)
library(car)
data(Prestige)
help(Prestige)
\end{verbatim} 


\noindent We would like to study whether individuals with higher levels of income have more prestigious jobs. Moreover, we would like to study whether professionals have more prestigious jobs than blue and white collar workers.

\newpage
\begin{enumerate}
	
	\item [(a)]
	\textbf{Create a new variable \texttt{professional} by recoding the variable \texttt{type} so that professionals are coded as $1$, and blue and white collar workers are coded as $0$ (Hint: \texttt{ifelse}).} \\
		
		\lstinputlisting[language=R, firstline=36, lastline=38]{PS04_answers_EGL.R}
			
		\begin{verbatim}
	  bc  prof  wc  <NA> 
			44   31   23   4 
		\end{verbatim}
			
			We can quickly check that there are 31 professionals, 44 blue collar workers, 23 white collar workers, and 4 occupations with the type \texttt{NA} in our dataset. \\
			
		\lstinputlisting[language=R, firstline=46, lastline=47]{PS04_answers_EGL.R}
			\begin{verbatim}
			0    1    <NA> 
			67   31    4 
			\end{verbatim}
			
			Now we can check that the 31 professionals were successfully recoded as \texttt{1} and the 67 blue and white collar workers as \texttt{0}.
			
	\vspace{1cm}
	
	
	\item [(b)]
	\textbf{Run a linear model with \texttt{prestige} as an outcome and \texttt{income}, \texttt{professional}, and the interaction of the two as predictors (Note: this is a continuous $\times$ dummy interaction.)} \\
		
				\lstinputlisting[language=R, firstline=50, lastline=53]{PS04_answers_EGL.R}
			
					\begin{table}[H] \centering 
			\caption{} 
			\label{} 
			\begin{tabular}{@{\extracolsep{5pt}}lc} 
				\\[-1.8ex]\hline 
				\hline \\[-1.8ex] 
				& \multicolumn{1}{c}{\textit{Dependent variable:}} \\ 
				\cline{2-2} 
				\\[-1.8ex] & prestige \\ 
				\hline \\[-1.8ex] 
				income & 0.003$^{***}$ \\ 
				& (0.0005) \\ 
				& \\ 
				professional & 37.781$^{***}$ \\ 
				& (4.248) \\ 
				& \\ 
				income:professional & $-$0.002$^{***}$ \\ 
				& (0.001) \\ 
				& \\ 
				Constant & 21.142$^{***}$ \\ 
				& (2.804) \\ 
				& \\ 
				\hline \\[-1.8ex] 
				Observations & 98 \\ 
				R$^{2}$ & 0.787 \\ 
				Adjusted R$^{2}$ & 0.780 \\ 
				Residual Std. Error & 8.012 (df = 94) \\ 
				F Statistic & 115.878$^{***}$ (df = 3; 94) \\ 
				\hline 
				\hline \\[-1.8ex] 
				\textit{Note:}  & \multicolumn{1}{r}{$^{*}$p$<$0.1; $^{**}$p$<$0.05; $^{***}$p$<$0.01} \\ 
			\end{tabular} 
		\end{table} 
		
		\begin{verbatim}
			Call:
			lm(formula = prestige ~ income + professional + income:professional, 
			data = Prestige)
			
			Residuals:
			Min      1Q  Median      3Q     Max 
			-14.852  -5.332  -1.272   4.658  29.932 
			
			Coefficients:
			Estimate Std. Error t value Pr(>|t|)    
			(Intercept)         21.1422589  2.8044261   7.539 2.93e-11 ***
			income               0.0031709  0.0004993   6.351 7.55e-09 ***
			professional        37.7812800  4.2482744   8.893 4.14e-14 ***
			income:professional -0.0023257  0.0005675  -4.098 8.83e-05 ***
			---
			Signif. codes:  0 ‘***’ 0.001 ‘**’ 0.01 ‘*’ 0.05 ‘.’ 0.1 ‘ ’ 1
			
			Residual standard error: 8.012 on 94 degrees of freedom
			(4 observations deleted due to missingness)
			Multiple R-squared:  0.7872,	Adjusted R-squared:  0.7804 
			F-statistic: 115.9 on 3 and 94 DF,  p-value: < 2.2e-16
		\end{verbatim}
		
		
	\item [(c)]
	\textbf{Write the prediction equation based on the result.} 
	
	\begin{align*}
	\widehat{prestige} &= \hat{\beta_0} + \hat{\beta_1} \times income + \hat{\beta_2} \times professional + \hat{\beta_3} \times (income \times professional)\\
	\widehat{prestige} &= 21.1422589 + 0.0031709 \times income + 37.78128 \times prof -0.0023257 \times (income \times prof)
	\end{align*}
	\vspace{.1cm}	
	
	\item [(d)]
	\textbf{Interpret the coefficient for \texttt{income}.} \\
		
		\noindent Because the model includes an interaction, $\hat{\beta_1}=0.0031709$ is a conditional effect, not an unconditional one. It is the effect of income for the baseline group, i.e.\ blue and white collar workers (\texttt{professional} $=0$): among non-professionals, a one-dollar increase in income is associated with an average increase of 0.003 points in prestige (equivalently, about 3.17 prestige points per \$1,000).
		
	\vspace{.5cm}	
	
	\item [(e)]
	\textbf{Interpret the coefficient for \texttt{professional}.} \\
		
		\noindent $\hat{\beta_2}=37.781$ is also a conditional effect. It is the average difference in prestige between professionals and non-professionals when \texttt{income} $=0$. Since no occupation has an income of zero, this should more be read as the shift in the intercept between the two regression lines rather than as a substantively meaningful quantity.
		
	\vspace{1cm}	
	
	\item [(f)]
	\textbf{What is the effect of a \$1,000 increase in income on prestige score for professional occupations? In other words, we are interested in the marginal effect of income when the variable \texttt{professional} takes the value of $1$. Calculate the change in $\hat{y}$ associated with a \$1,000 increase in income based on your answer for (c).}
	\vspace{.5cm}
		
			\lstinputlisting[language=R, firstline=56, lastline=59]{PS04_answers_EGL.R}
			\lstinputlisting[language=R, firstline=72, lastline=72]{PS04_answers_EGL.R}
			
			\begin{verbatim}
				0.8452
			\end{verbatim}
		
		The marginal effect of income is $\hat{\beta_1}+\hat{\beta_3}$:
			\begin{align*}
				\Delta \widehat{prestige} &= (0.0031709 - 0.0023257) \times 1000 \approx 0.845
			\end{align*}
			
		For professionals, a \$1,000 increase in income is associated with an average increase of about 0.85 prestige points, much smaller than the 3.17 points for non-professionals.
		
	\vspace{1cm}
	\item [(g)]
	\textbf{What is the effect of changing one's occupations from non-professional to professional when her income is \$6,000? We are interested in the marginal effect of professional jobs when the variable \texttt{income} takes the value of $6,000$. Calculate the change in $\hat{y}$ based on your answer for (c).} \\
		
		\lstinputlisting[language=R, firstline=75, lastline=75]{PS04_answers_EGL.R}
		
		\begin{verbatim}
			23.82703
		\end{verbatim}
		
		The marginal effect of \texttt{professional} is $\hat{\beta_2}+\hat{\beta_3}\times income$:
		\begin{align*}
			\Delta \widehat{prestige} &= 37.781 + (-0.0023257 \times 6000) \approx 23.83
		\end{align*}
		
		At an income of \$6,000, moving from a non-professional to a professional occupation is associated with an average increase of about 23.83 prestige points.
	
	
\end{enumerate}

\newpage

\section*{Question 2: Political Science}
\vspace{.25cm}
\noindent 	Researchers are interested in learning the effect of all of those yard signs on voting preferences.\footnote{Donald P. Green, Jonathan	S. Krasno, Alexander Coppock, Benjamin D. Farrer,	Brandon Lenoir, Joshua N. Zingher. 2016. ``The effects of lawn signs on vote outcomes: Results from four randomized field experiments.'' Electoral Studies 41: 143-150. } Working with a campaign in Fairfax County, Virginia, 131 precincts were randomly divided into a treatment and control group. In 30 precincts, signs were posted around the precinct that read, ``For Sale: Terry McAuliffe. Don't Sellout Virgina on November 5.'' \\

Below is the result of a regression with two variables and a constant.  The dependent variable is the proportion of the vote that went to McAuliff's opponent Ken Cuccinelli. The first variable indicates whether a precinct was randomly assigned to have the sign against McAuliffe posted. The second variable indicates
a precinct that was adjacent to a precinct in the treatment group (since people in those precincts might be exposed to the signs).  \\

\vspace{.5cm}
\begin{table}[!htbp]
	\centering 
	\textbf{Impact of lawn signs on vote share}\\
	\begin{tabular}{@{\extracolsep{5pt}}lccc} 
		\\[-1.8ex] 
		\hline \\[-1.8ex]
		Precinct assigned lawn signs  (n=30)  & 0.042\\
		& (0.016) \\
		Precinct adjacent to lawn signs (n=76) & 0.042 \\
		&  (0.013) \\
		Constant  & 0.302\\
		& (0.011)
		\\
		\hline \\
	\end{tabular}\\
	\footnotesize{\textit{Notes:} $R^2$=0.094, N=131}
\end{table}

\vspace{1cm}
\noindent For all tests below, the degrees of freedom are $n-k-1 = 131-2-1 = 128$.

\lstinputlisting[language=R, firstline=80, lastline=82]{PS04_answers_EGL.R}

\vspace{.5cm}
\begin{enumerate}
	\item [(a)] \textbf{Use the results from a linear regression to determine whether having these yard signs in a precinct affects vote share (e.g., conduct a hypothesis test with $\alpha = .05$).}
		
		\lstinputlisting[language=R, firstline=87, lastline=89]{PS04_answers_EGL.R}
		
			\noindent \textbf{Step 1 (assumptions):} Precincts were randomly assigned to treatment, and the OLS assumptions (linearity, zero conditional mean, constant variance, normal errors) are taken to hold.\\
			\textbf{Step 2 (hypotheses):} $H_0: \beta_1 = 0$ versus $H_A: \beta_1 \neq 0$.\\
			\textbf{Step 3 (test statistic):} $t = \dfrac{\hat{\beta_1}-0}{SE(\hat{\beta_1})} = \dfrac{0.042}{0.016} = 2.625$\\
			\textbf{Step 4 (p-value):} $p = 2 \times Pr(t_{128} > |2.625|) \approx 0.0097$\\
			\textbf{Step 5 (conclusion):} Since $p < 0.05$, we reject $H_0$. \\
			
		\lstinputlisting[language=R, firstline=90, lastline=90]{PS04_answers_EGL.R}
		\begin{verbatim}
			0.01034127 0.07365873
		\end{verbatim}
			
			\noindent The 95\% confidence interval is $[0.010, 0.074]$, which excludes zero. There is a statistically reliable relationship: precincts assigned lawn signs had, on average, a 4.2 percentage point higher vote share for Cuccinelli, the opponent of the candidate named on the signs, than control precincts (holding adjacency constant). 
			
			Because precincts were randomly assigned to treatment, the results support a causal interpretation of the effect of lawn-sign assignment, although potential spillovers between precincts complicate the interpretation.
		
	\vspace{1cm}
	\item [(b)]  \textbf{Use the results to determine whether being next to precincts with these yard signs affects vote share (e.g., conduct a hypothesis test with $\alpha = .05$).} \\
		
		\lstinputlisting[language=R, firstline=98, lastline=100]{PS04_answers_EGL.R}
			\noindent \textbf{Hypotheses:} $H_0: \beta_2 = 0$ versus $H_A: \beta_2 \neq 0$.\\
			\textbf{Test statistic:} $t = \dfrac{0.042-0}{0.013} \approx 3.231$\\
			\textbf{p-value:} $p = 2 \times Pr(t_{128} > |3.231|) \approx 0.0016$\\
			\textbf{Conclusion:} Since $p < 0.05$, we reject $H_0$. \\
			
		\lstinputlisting[language=R, firstline=101, lastline=101]{PS04_answers_EGL.R}
		\begin{verbatim}
			0.01627728 0.06772272
		\end{verbatim}
			
			\noindent The 95\% confidence interval is $[0.016, 0.068]$, which excludes zero. Precincts adjacent to treated precincts also had, on average, a 4.2 percentage point higher vote share for Cuccinelli than precincts that were neither treated nor adjacent. This suggests the signs' effect spills over to neighbouring precincts, so the treatment effect is not confined to the precincts in which the signs were posted.
		
	\vspace{1cm}
	\item [(c)] \textbf{Interpret the coefficient for the constant term substantively.} \\
		
		\noindent $\hat{\beta_0} = 0.302$ is the predicted value of the outcome when all predictors equal zero. Here, that is a precinct that was neither assigned lawn signs nor adjacent to a precinct with signs (the control group). The expected vote share for Cuccinelli here is 30.2\%. The constant is highly distinguishable from zero ($t = 0.302/0.011 \approx 27.5$). \\
		
		\lstinputlisting[language=R, firstline=112, lastline=112]{PS04_answers_EGL.R}
		\begin{verbatim}
			27.45455
		\end{verbatim}
		
	\vspace{.5cm}
	\item [(d)] \textbf{Evaluate the model fit for this regression. What does this	tell us about the importance of yard signs versus other factors that are not modeled?} \\
		
		\noindent $R^2 = 0.094$, which suggests the two predictors together explain only 9.4\% of the variation in Cuccinelli's vote share across precincts. About 90.6\% of the variation is left unexplained. The model fits poorly in terms of predictive power—although the yard signs have statistically significant effects, the model explains only a small proportion of the overall variation in vote share. This does not necessarily mean that the effects of yard signs are substantively unimportant: a 4.2-percentage-point increase could be consequential in a close election. Rather, the low \(R^2\) suggests that other factors, such as partisanship and precinct demographics, may account for much of the variation in electoral outcomes.
		
\end{enumerate}  
\end{document}
